\subsection{The presentable case}
\label{subsec:2dloc-presentable}

\gutentag{0081}%
In this subsection $\D \in \tprl$ is a presentable $(\infty,2)$-category and $I$, $P \subset \func(\Delta^{1},\D)$ are \emph{essentially small} locally full sub-$(\infty,2)$-categories, e.g.\ generated by sets of $1$-morphisms and commutative squares. Presentability provides the left adjoint to $\iota\colon \Pc_{I,P} \to \D$, and with it the identification of $\Pc_{I,P}$ with a pushout.
For large $I$, $P$ the results of \cref{subsec:2dloc-local,subsec:2dloc-plain,subsec:2dloc-saturation} still apply, but the argument below will not produce a left adjoint.

\subsubsection*{Presentable $(\infty,2)$-categories}
\gutentag{0082}%
A \emph{presentable $(\infty,2)$-category} is a $\Cat$-module in $\prl$; we write $\tprl := \modcat{\Cat}{\prl}$ for the $\infty$-category of presentable $\infty$-categories with a colimit-preserving $\Cat$-action (for the cartesian product) and $\Cat$-linear left adjoints.
By \cite{gepnerhaugseng2015enriched,heine2023enriched} every $\D \in \tprl$ is canonically $\Cat$-enriched, with $\func(K,\Hom_{\D}(d,d')) \simeq \mapp_{\D}(K \otimes d,d')$, tensored and cotensored over $\Cat$, and $\Cat$-linear functors are canonically enriched.
\todo{Precise references.}
For small $\A$, $\Free(\A) := \func(\A\opcat,\Cat)$ is in $\tprl$, and by the enriched Yoneda lemma and the universal property of enriched presheaf categories \cite{hinich2020yoneda,heine2023enriched}, restriction along the Yoneda embedding $y\colon \A \to \Free(\A)$ is an equivalence
\begin{equation}
\label{eq:free-univ}
\func^{L}_{\Cat}\bigl(\Free(\A),\D\bigr) \xrightarrow{\ \simeq\ } \func(\A,\D)
\end{equation}
for every $\D \in \tprl$, where the left-hand side is the $(\infty,2)$-category of $\Cat$-linear left adjoints.
A functor $f\colon \A \to \B$ induces $\Free(f) = f_{!}$ (enriched left Kan extension), left adjoint to restriction $f^{*}$, and $\Free\colon \twoCat \to \tprl$ preserves colimits.
Let $\Edge := \Free(\Delta^{1})$. By \eqref{eq:free-univ} a $1$-morphism $s\colon x \to y$ of $\D \in \tprl$ is a map $s\colon \Edge \to \D$ in $\tprl$, and by \cite[Thm.~4.4]{haugseng2021laxmonads} a left adjoint $s$ extends essentially uniquely along $\Free(\ell)$ to $\Free(\Adj) \to \D$.

\subsubsection*{Colimits}
\gutentag{0083}%
Colimits in $\iota_{1}\D$ are colimits in the $(\infty,2)$-categorical sense: $\Hom_{\D}(-,e)$ turns them into limits of $\infty$-categories, since $\func(K,\Hom_{\D}(-,e))^{\simeq} \simeq \mapp_{\D}(K \otimes -,e)$ for $K \in \Cat$, $K \otimes -$ preserves colimits, and the Yoneda lemma applies in $K$. In particular all $(\infty,2)$-pushouts exist, and every $s\colon x \to y$ has the cylinder and cocylinder
\[
\Cyl(s) = (\Delta^{1} \otimes x) \sqcup_{\{1\} \otimes x,\,s} y \ , \qquad \Cyl'(s) = y \sqcup_{s,\,\{0\} \otimes x} (\Delta^{1} \otimes x) \ ,
\]
with $j_{s}\colon x \simeq \{0\} \otimes x \to \Cyl(s)$ and $k_{s}\colon x \simeq \{1\} \otimes x \to \Cyl'(s)$: as $\Hom_{\D}(\Delta^{1} \otimes x,e) \simeq \func(\Delta^{1},\Hom_{\D}(x,e))$, we have
\[
\Hom_{\D}(\Cyl(s),e) \simeq \func(\Delta^{1},\Hom_{\D}(x,e)) \times_{\Hom_{\D}(x,e)} \Hom_{\D}(y,e) \simeq (\mathrm{id} \downarrow s^{\natural}) \ ,
\]
and likewise for $\Cyl'(s)$.

\subsubsection*{The localization as a pushout}

\begin{defi}
\label{def:adjoin-pres}
\gutentag{0084}%
By \eqref{eq:free-univ}, the evaluations $\Delta^{1} \times I \to \D$ and $\Delta^{1} \times P \to \D$ are maps $\Free(\Delta^{1} \times I) \to \D$ and $\Free(\Delta^{1} \times P) \to \D$ in $\tprl$. The pushout in $\tprl$
\[
\begin{tikzcd}[column sep=large]
\Free(\Delta^{1} \times I) \oplus \Free(\Delta^{1} \times P) \ar[r, "\mathrm{ev}"] \ar[d, "\Free(\rho \times \mathrm{id}) \oplus \Free(\ell \times \mathrm{id})"'] & \D \ar[d, "L"] \\
\Free(\Refl \times I) \oplus \Free(\Corefl \times P) \ar[r] & \D\lrrefl{I}{P}
\end{tikzcd}
\]
is again denoted $\D\lrrefl{I}{P}$; likewise $\D\lradj{I}{P}$ with $\Adj$ in place of $\Refl$ and $\Corefl$.
\end{defi}

\gutentag{0085}%
As $\Free$ preserves colimits, $\Free(\A\lrrefl{I}{P}) \simeq \Free(\A)\lrrefl{y(I)}{y(P)}$ for small $\A$, with $\A\lrrefl{I}{P}$ as in \cref{def:adjoin} and $y(I) \subset \func(\Delta^{1},\Free(\A))$ generated by the image of $I$ (as in the paragraph after \cref{def:adjoin}, only the generated locally full sub-$(\infty,2)$-category matters).
By \eqref{eq:free-univ}, $\func^{L}_{\Cat}(\D\lrrefl{I}{P},\E)$ is the pullback of $\func^{L}_{\Cat}(\D,\E) \to \func(I,\func(\Delta^{1},\E)) \times \func(P,\func(\Delta^{1},\E))$ along $(\rho^{*})_{*} \times (\ell^{*})_{*}$, so the proof of \cref{prop:adjoin-univ} gives:

\begin{prop}
\label{prop:adjoin-univ-pres}
\gutentag{0086}%
For $\E \in \tprl$, restriction along $L$ identifies $\func^{L}_{\Cat}(\D\lrrefl{I}{P},\E)$ with the locally full sub-$(\infty,2)$-category of $\func^{L}_{\Cat}(\D,\E)$ described in \cref{prop:adjoin-univ}.
\end{prop}

\begin{thm}
\label{thm:2dloc-local}
\gutentag{0087}%
Let $\D \in \tprl$ and let $I$, $P \subset \func(\Delta^{1},\D)$ be essentially small locally full sub-$(\infty,2)$-categories.
\begin{enumerate}
\item $L\colon \D \to \D\lrrefl{I}{P}$ has a right adjoint, which is a locally full inclusion with image $\Pc_{I,P}$. Hence $\D\lrrefl{I}{P} \simeq \Pc_{I,P}$, and $\iota\colon \Pc_{I,P} \to \D$ has the left adjoint $L$ in the sense of \cref{def:enriched-adjoint}.
\item $T := \iota L$ is a $\Cat$-enriched monad on $\D$ and $\iota$ is monadic, so $\Pc_{I,P} \simeq \Alg_{T}(\D)$. The space of $T$-algebra structures on $d \in \D$ is empty or contractible, and non-empty if and only if $d$ is $(I,P)$-local.
\end{enumerate}
\end{thm}

\begin{proof}
\gutentag{0088}%
(1) \emph{The underlying $\infty$-category.}
Colimits in $\tprl = \modcat{\Cat}{\prl}$ are computed in $\prl$ \cite[4.2.3.5]{ha}\todo{check the tag}, and pushouts in $\prl$ as pullbacks of right adjoints in $\widehat{\Cat}$ \cite[5.5.3.4, 5.5.3.18]{htt}. The right adjoints are:
\begin{itemize}
\item the right adjoint of $\Free(\Delta^{1} \times I) \to \D$. By the enriched Yoneda lemma \cite{hinich2020yoneda,heine2023enriched}, it sends $d$ to $\Hom_{\D}(\mathrm{ev}(-),d)$, so it is the underlying functor of $R_{I}\colon \D \to \func(I\opcat,\func(\Delta^{1},\Cat))$ of \cref{subsec:2dloc-local}, under $\func((\Delta^{1})\opcat,\Cat) \simeq \func(\Delta^{1},\Cat)$; likewise for $P$;
\item the right adjoint of $\Free(\ell \times \mathrm{id}) = (\ell \times \mathrm{id})_{!}$, restriction along $(\ell \times \mathrm{id})\opcat$. The equivalence $\Adj \simeq \Adj\opcat$, $l \mapsto r\opcat$, $r \mapsto l\opcat$ (in $\Adj\opcat$, $r\opcat \dashv l\opcat$ with the same unit and counit) carries $\rho$ to $\ell\opcat$ and preserves the unit, so it induces $\Corefl \simeq \Corefl\opcat$ and this right adjoint is $(\rho^{*})_{*}\colon \func(P\opcat,\func(\Corefl,\Cat)) \to \func(P\opcat,\func(\Delta^{1},\Cat))$. Likewise the right adjoint of $\Free(\rho \times \mathrm{id})$ is $(\ell^{*})_{*}\colon \func(I\opcat,\func(\Refl,\Cat)) \to \func(I\opcat,\func(\Delta^{1},\Cat))$.
\end{itemize}
So $\iota_{1}\D\lrrefl{I}{P}$ is the pullback of the underlying $\infty$-categories of \eqref{eq:pullback}, with projection to $\iota_{1}\D$ the right adjoint $\iota$ of $L$. As $\iota_{1}$ preserves limits, \cref{thm:local-pullback} identifies $\iota$ with the inclusion $\iota_{1}\Pc_{I,P} \subset \iota_{1}\D$.

\emph{The enrichment.}
By \cite[7.3.2.7]{ha}, $\iota$ is canonically lax $\Cat$-linear, hence enriched \cite{heine2023enriched}.
\todo{Check that the passage from lax-linear functors and transformations to enriched ones in \cite{heine2023enriched} is compatible with natural transformations.}
As $L$ is $\Cat$-linear, $\iota$ preserves cotensors: naturally in $d$,
\[
\mapp_{\D}(d,\iota(K \pitchfork q)) \simeq \mapp(K \otimes Ld,q) \simeq \mapp(L(K \otimes d),q) \simeq \mapp_{\D}(d,K \pitchfork \iota q) \ .
\]
For $q,q' \in \D\lrrefl{I}{P}$ and $K \in \Cat$ we get
\[
\func\bigl(K,\Hom_{\D\lrrefl{I}{P}}(q,q')\bigr)^{\simeq} \simeq \mapp(q,K \pitchfork q') \subset \mapp_{\D}(\iota q,K \pitchfork \iota q') \simeq \func\bigl(K,\Hom_{\D}(\iota q,\iota q')\bigr)^{\simeq} \ ,
\]
where by the first part the left-hand side consists of the $(I,P)$-local $g\colon \iota q \to K \pitchfork \iota q'$.
Now
\[
\Hom_{\D}(s,K \pitchfork \iota q') = \func(K,\Hom_{\D}(s,\iota q'))
\]
has the fully faithful right adjoint $\func(K,(i^{\natural})^{R})$ for an object $s = i$ of $I$ and the fully faithful left adjoint $\func(K,(p^{\natural})^{L})$ for an object $s = p$ of $P$, and evaluation at $k \in K$ commutes with $s^{\natural}$, these adjoints, units and counits. So the mates for $g$ evaluate at $k$ to the mates for $g_{k}\colon \iota q \to \iota q'$, and $g$ is $(I,P)$-local if and only if every $g_{k}$ is, i.e.\ if and only if the corresponding functor $K \to \Hom_{\D}(\iota q,\iota q')$ factors through the full subcategory of $(I,P)$-local $1$-morphisms. By Yoneda in $K$, $\Hom_{\D\lrrefl{I}{P}}(q,q') \to \Hom_{\D}(\iota q,\iota q')$ is the inclusion of this full subcategory, i.e.\ agrees with $\Hom_{\Pc_{I,P}}(\iota q,\iota q') \subset \Hom_{\D}(\iota q,\iota q')$.
So the enriched $\iota$ is locally fully faithful and, by the first part for $K = \Delta^{0}$, a monomorphism on objects, i.e.\ a locally full inclusion with image $\Pc_{I,P}$ (\cref{subsec:prelim-conventions}). By \cite{heine2023enriched}, $L \dashv \iota$ is an adjunction in the $(\infty,2)$-category $\twoCat$ of large $(\infty,2)$-categories, with enriched unit and counit, as required in \cref{def:enriched-adjoint}.

(2) Hence $T = \iota L$ is enriched, with enriched $\eta$ and $\mu = \iota\varepsilon L$ satisfying the monad identities in the homotopy $2$-category. On underlying $\infty$-categories $T$ is the monad of $L \dashv \iota$, an associative algebra in $\func(\iota_{1}\D,\iota_{1}\D)$ \cite[\S 4.7.3]{ha}; monads in an $(\infty,2)$-category are compared with such algebras in \cite{haugseng2021laxmonads}.
We check the hypotheses of Barr--Beck--Lurie \cite[4.7.3.5]{ha} for $\iota$ on underlying $\infty$-categories.
$\iota$ is conservative: if $f$ is $(I,P)$-local and invertible in $\D$, then by \cref{lem:mate-correspondence}(3) the mates for $f^{-1}$ are inverse to those for $f$, so $f^{-1}$ is $(I,P)$-local.
$\iota$ preserves colimits of $\iota$-split simplicial objects: let $q_{\bullet}$ be a simplicial object of $\Pc_{I,P}$ with $\iota q_{\bullet}$ split, and write $q_{\bullet} = (\iota q_{\bullet},F_{\bullet})$ in \eqref{eq:pullback}. Colimits of split simplicial objects are preserved by every functor \cite[6.1.3.16]{htt}, so $|\iota q_{\bullet}|$ exists and is preserved by $(R_{I},R_{P})$; and $|F_{\bullet}|$ exists in $\func(I\opcat,\func(\Refl,\Cat)) \times \func(P\opcat,\func(\Corefl,\Cat))$ and is preserved by $(\ell^{*})_{*} \times (\rho^{*})_{*}$, colimits being pointwise. So $(|\iota q_{\bullet}|,|F_{\bullet}|)$ is a colimit of $q_{\bullet}$ preserved by $\iota$.
Thus $\iota$ is monadic, $\Pc_{I,P} \simeq \Alg_{T}(\D)$ over $\D$, and as $\iota$ is a monomorphism the fibres of $\Alg_{T}(\D) \to \D$ are empty or contractible, non-empty exactly over the $(I,P)$-local objects.
\end{proof}

\gutentag{0089}%
The proof with $\Adj$ in place of $\Refl$ and $\Corefl$, and the pullback description of $\Pc^{\dashv}_{I,P}$ in \cref{subsec:2dloc-plain}, give $\D\lradj{I}{P} \simeq \Pc^{\dashv}_{I,P}$ with the analogous monadicity statement.
By \cref{prop:adjoin-univ-pres}, $L^{*}\colon \func^{L}_{\Cat}(\D\lrrefl{I}{P},\E) \to \func^{L}_{\Cat}(\D,\E)$ is a locally full inclusion for every $\E \in \tprl$.

\subsubsection*{Consequences of the plain results}
\gutentag{008A}%
As $\iota$ has a left adjoint, the results of \cref{subsec:2dloc-local,subsec:2dloc-plain,subsec:2dloc-saturation} specialize as follows.

\begin{cor}
\label{cor:presentable}
\gutentag{008B}%
Let $\D \in \tprl$ and $I$, $P \subset \func(\Delta^{1},\D)$ essentially small locally full sub-$(\infty,2)$-categories.
\begin{enumerate}
\item $Li$ is reflective, $L\sigma$ left Beck--Chevalley, $Lp$ coreflective and $L\tau$ right Beck--Chevalley in $\D\lrrefl{I}{P}$, for the objects $i$, $p$ and $1$-morphisms $\sigma$, $\tau$ of $I$ and $P$; and $(i^{\natural})^{R} \simeq ((Li)^{L})^{\natural}$, $(p^{\natural})^{L} \simeq ((Lp)^{R})^{\natural}$ (\cref{prop:local-adjoint-refl}).
\item $\D\refl{I}\corefl{L_{I}(P)} \simeq \D\lrrefl{I}{P} \simeq \D\corefl{P}\refl{L_{P}(I)}$ as locally full sub-$(\infty,2)$-categories of $\D$ (\cref{prop:2dloc-commute}), where $L_{I}\colon \D \to \D\refl{I}$ and $L_{P}\colon \D \to \D\corefl{P}$, and $L_{I}(P)$, $L_{P}(I)$ are generated by the images of $P$ and $I$.
\item (\cref{prop:cylinder}.) We have
\[
\D\lradj{I}{P} \simeq \Pc^{\dashv}_{I,P} = \Pc_{K_{I},J_{P}} \simeq \D\lrrefl{K_{I}}{J_{P}} \ .
\]
\item (\cref{prop:saturation-L}.) A $1$-morphism $f$ is in $\Sat^{\mathrm{corefl}}_{\Pc_{I,P}}$, resp.\ $\Sat^{\mathrm{refl}}_{\Pc_{I,P}}$, if and only if $Lf$ is coreflective, resp.\ reflective, in $\D\lrrefl{I}{P}$, and a square $\sigma$ between such $1$-morphisms is a square of the saturation if and only if $L\sigma$ is right, resp.\ left, Beck--Chevalley. In particular $\overline{P}$ consists of the $f$ with $Lf$ coreflective in $\D\corefl{P}$ and the $\sigma$ with $L\sigma$ right Beck--Chevalley, and $\overline{P}^{\dashv}$ of the $f$ with $Lf$ a left adjoint in $\D\radj{P}$ and the $\sigma$ with $L\sigma$ right Beck--Chevalley.
\item (\cref{thm:sat-units}.) If every unit $\eta_{d}\colon d \to \iota L d$ of $\D \to \D\corefl{P}$ lies in the coreflectively saturated class $\langle P \rangle$ generated by $P$, then $\langle P \rangle$ and $\overline{P}$ have the same objects.
\end{enumerate}
\end{cor}

\begin{proof}
\gutentag{008C}%
(1)--(3) are the cited results together with \cref{thm:2dloc-local}; in (2) the identifications are unique over $\D$, since a locally full sub-$(\infty,2)$-category is determined by its objects and $1$-morphisms, so they agree with the one from pasting pushouts.
(4) and (5) are \cref{prop:saturation-L}(2) and \cref{thm:sat-units} with \cref{thm:2dloc-local}; the description of $\overline{P}$ is \cref{lem:saturation-explicit}. For $\overline{P}^{\dashv}$: by the plain version of \cref{lem:repr-refl} applied in $\Pc^{\dashv}_{\emptyset,P}$ (\cref{subsec:2dloc-plain}), $Lf$ is a left adjoint if and only if every object and $1$-morphism of $\Pc^{\dashv}_{\emptyset,P}$ is $(\emptyset,\{Lf\})$-local in the plain sense, and under $\Hom_{\Pc^{\dashv}_{\emptyset,P}}(L-,q) \simeq \Hom_{\D}(-,\iota q)$ this says that $f$ is in $\overline{P}^{\dashv}$; likewise for squares.
\end{proof}

\begin{rmk}
\label{rmk:free-local}
\gutentag{008D}%
For small $\A$ and small locally full $I$, $P \subset \func(\Delta^{1},\A)$, \cref{thm:2dloc-local} gives
\[
\Free(\A\lrrefl{I}{P}) \simeq \Free(\A)\lrrefl{y(I)}{y(P)} \simeq \Pc_{y(I),y(P)}(\Free(\A)) \ .
\]
By the Yoneda lemma $\Hom(y(s),X) \simeq X(s)$, so a presheaf $X\colon \A\opcat \to \Cat$ is $(y(I),y(P))$-local if and only if it satisfies \cref{thm:local-pullback}(1) with $X(s)$ in place of $s^{\natural}$: every $X(i)$ has a fully faithful right adjoint, every $X(p)$ a fully faithful left adjoint, and $X$ takes the $1$-morphisms of $I$ and $P$ to right, resp.\ left, Beck--Chevalley squares. By the paragraph after \cref{thm:local-pullback}, these are the $X$ that extend along $\gamma\opcat$, in accordance with the composite equivalence, which is restriction along $\gamma\opcat$.
\end{rmk}

\begin{rmk}
\label{rmk:eilenberg-moore-enriched}
\gutentag{008E}%
\Cref{thm:2dloc-local}(2) identifies $\Pc_{I,P}$ with $\Alg_{T}(\D)$ as $\infty$-categories. On $2$-morphisms: if $f,g\colon \iota q \to \iota q'$ are $(I,P)$-local, every $\alpha\colon f \Rightarrow g$ is an algebra $2$-morphism, i.e.\ $a_{q'}T(\alpha) = \alpha a_{q}$ as $2$-morphisms $a_{q'}T(f) \simeq fa_{q} \Rightarrow ga_{q} \simeq a_{q'}T(g)$, where $a_{q} = \iota\varepsilon_{q}$. Indeed, both sides are $2$-morphisms between $(I,P)$-local $1$-morphisms $T\iota q \to \iota q'$ restricting along $\eta_{\iota q}$ to $\alpha$, and restriction along $\eta_{\iota q}$ is an equivalence $\Hom_{\Pc_{I,P}}(L\iota q,q') \simeq \Hom_{\D}(\iota q,\iota q')$ by the adjunction, and $\Pc_{I,P}$ is locally full. This is a computation in $\htwo\D$. So $T$-algebras, strong morphisms and algebra $2$-morphisms form an $(\infty,2)$-category with the same hom-categories as $\Pc_{I,P}$.
\todo{Compare with the Eilenberg--Moore object of \cite{riehlverity2016adjunctions,haugseng2021laxmonads}.}
\end{rmk}

\begin{rmk}
\label{rmk:2dloc-compare-bousfield}
\gutentag{008F}%
In \cref{prop:bousfield-pushout} the local objects are those $c$ for which $s^{\natural}$ is an \emph{equivalence}; here $i^{\natural}$ is a \emph{localization} and $p^{\natural}$ a \emph{colocalization}.
The Beck--Chevalley condition on $1$-morphisms has no one-dimensional counterpart, since, unlike adjoints, inverses are preserved by every square.
It is what prevents $\iota$ from being fully faithful and $T$ from being idempotent: in general $L\iota q \not\simeq q$, because $\Hom_{\Pc_{I,P}}(L\iota q,q') \simeq \Hom_{\D}(\iota q,\iota q')$ also contains the non-local $1$-morphisms.
\end{rmk}

% The subsection on lax-idempotence of T (and the remarks on two-sided adjoining and six-functor formalisms)
% was removed; it is in commit 31c3cb5625af320b08fe8cb95ade64f62acf416a, file sections/2dloc/06-laxidempotent.tex.
